% !TeX root = ../informe.tex
% ===========================================================================
%  7. Resultados
%
%  Cierra el círculo con los objetivos (sección 2) y con los criterios de
%  aceptación acordados. Los \label{ca:N} son los que hacen funcionar
%  \crit{N} desde el resto del documento: si añades filas, añade su \label.
% ===========================================================================

\section{Resultados}
\label{sec:resultados}

\subsection{Funcionalidades implementadas}
\label{sec:funcionalidades}

\instruccion{Qué quedó operativo, con capturas de las pantallas principales.
Dos o tres capturas bien elegidas valen más que ocho: las que muestran el
flujo central del sistema. El resto va al anexo de capturas.}

\ph{Resumen de lo implementado.}

\subsection{Cumplimiento de los criterios de aceptación}
\label{sec:cumplimiento}

\instruccion{Un criterio por fila, con la evidencia que respalda cada «Sí».
Es la tabla que un evaluador busca primero. Si alguno no se cumple,
decláralo: un criterio marcado como cumplido que luego falla en la demo
cuesta más que uno declarado pendiente.}

\begin{table}[H]
  \centering
  \caption{Cumplimiento de los criterios de aceptación.}
  \label{tab:cumplimiento}
  \begin{tabularx}{\textwidth}{C{1.4cm} Y C{2cm} L{3cm}}
    \toprule
    \textbf{\#} & \textbf{Criterio} & \textbf{Estado} & \textbf{Evidencia} \\
    \midrule
    \textsf{CA-1}\label{ca:1} & \ph{Criterio acordado} & \ph{...} & \ph{...} \\
    \textsf{CA-2}\label{ca:2} & \ph{Criterio acordado} & \ph{...} & \ph{...} \\
    \textsf{CA-3}\label{ca:3} & \ph{Criterio acordado} & \ph{...} & \ph{...} \\
    \bottomrule
  \end{tabularx}
\end{table}

\subsection{Verificación de las reglas de negocio}
\label{sec:verificacion-reglas}

\instruccion{Cómo se comprobó que cada regla de la \cref{tab:reglas}
funciona. No hace falta un capítulo de pruebas: una tabla con el escenario,
el resultado esperado y el obtenido es suficiente y se lee en treinta
segundos. La colección de peticiones va al anexo.}

\begin{table}[H]
  \centering
  \caption{Comprobación de las reglas de negocio.}
  \label{tab:pruebas-reglas}
  \begin{tabular}{C{1.6cm} L{6.6cm} L{4.8cm}}
    \toprule
    \textbf{Regla} & \textbf{Escenario} & \textbf{Resultado} \\
    \midrule
    \ph{...} & \ph{...} & \ph{...} \\
    \bottomrule
  \end{tabular}
\end{table}

\subsection{Limitaciones}
\label{sec:limitaciones}

\instruccion{Qué no se alcanzó y por qué. Declararlo suma: demuestra que se
conoce el propio sistema. Omitir una limitación que el evaluador descubre en
la demo resta el doble.}

\ph{Limitaciones conocidas.}
